% draft/VCH_CASESTUDY.tex
% VCH Case Study — Sections 6.2 (Data) and 6.3 (Results)
%
% Three models reported:
%   1. BIP (LP relaxation)   — static snapshot, theoretical minimum declinations
%   2. Queue Model           — dynamic online simulation (erlang_a / cyclic_fifo / optimised)
%   3. DES                   — extended simulation (geo + escalation + batch arrivals)
%
% NOTE: All model coefficients are fit to synthetic data mirroring VCH PARIS structure;
% framework is ready for replacement with real data once received from VCH.

\subsection{Data}
\label{sec:data}

\subsubsection*{System Parameters}

All parameters are calibrated against the 2024/25 BC Office of Seniors Advocate
Long-Term Care and Assisted Living Directory~\cite{ref9}.
Table~\ref{tab:vch_params} summarises the key VCH statistics and their translation
into model parameters.

\begin{table}[H]
\centering
\caption{VCH system parameters calibrated to OSA 2024/25.}
\label{tab:vch_params}
\begin{tabular}{llr}
\toprule
\textbf{Parameter} & \textbf{Source / derivation} & \textbf{Value} \\
\midrule
Number of facilities            & OSA 2024/25, p.\ 31              & 55 \\
\quad Health-authority-owned    &                                  & 16 \\
\quad Contracted (non-HA)       &                                  & 39 \\
Mean beds per facility          & $28{,}869 / 301$ (BC average)   & 96 \\
Estimated VCH beds              & $55 \times 96$                  & $\approx 5{,}300$ \\
Mean length of stay             & OSA 2024/25, p.\ 16             & 1{,}083 days \\
\midrule
Total vacancy rate              & $5{,}300 / 1{,}083$             & 147/month \\
Waitlist-accessible vacancies $\mu$ & $\approx 12\%$ to priority list & 18/month \\
Arrival rate $\lambda$          & Steady-state calibration        & 28/month \\
Initial waitlist backlog        & VCH reported, OSA 2024/25      & $\approx 500$ \\
\midrule
Mean wait --- all admissions      & OSA 2024/25, p.\ 14           & 81 days \\
Mean wait --- non-urgent community & OSA 2024/25, p.\ 14          & 473 days \\
\bottomrule
\end{tabular}
\end{table}

\subsubsection*{Priority Tier Fractions and Abandonment}

Four admission categories are modelled: Acute Care, Community Emergency,
Community High, and Transfer/Site Specific.  Tier fractions are calibrated so that
the Cyclic FIFO baseline reproduces the OSA-reported mean wait of 473 days for
non-urgent community (Community High) clients:
\[
  W_{\mathrm{CommHigh}} \approx
  \frac{\lambda_{\mathrm{CommHigh}} - \mu_{\mathrm{CommHigh}}}
       {\mu_{\mathrm{CommHigh}} \cdot \theta_{\mathrm{CommHigh}}} = 473\ \text{days.}
\]

\begin{table}[H]
\centering
\caption{Priority tier fractions and abandonment rates, calibrated to OSA 2024/25.}
\label{tab:tier_fracs}
\begin{tabular}{lrrl}
\toprule
\textbf{Tier} & \textbf{Fraction} & \textbf{$\theta$ (monthly)} & \textbf{Rationale} \\
\midrule
Acute Care             & 35.0\% & 50\% & Hospital cannot hold indefinitely ($\approx$2-month patience) \\
Community Emergency    & 30.0\% & 30\% & Moderate community support ($\approx$3-month patience) \\
Community High         & 21.1\% & 2\%  & Non-urgent; very patient; tuned to $W_{\text{CommHigh}} \approx 473$\,d \\
Transfer/Site Specific & 13.9\% & 15\% & Sending facility moderate pressure ($\approx$7-month patience) \\
\bottomrule
\end{tabular}
\end{table}

\subsubsection*{Facility Data}

Mock facility data (55 facilities) are generated to match the OSA 2024/25 distribution:
bed counts drawn from $\mathcal{N}(96, 28^2)$ (minimum 20), approximately 60\% of
facilities with a female gender restriction, and facility-specific random effects
$u_n \sim \mathcal{N}(0, 0.25)$.  Gender compatibility is enforced as a hard feasibility
constraint --- clients are not offered facilities whose gender restriction does not match.
All logistic model coefficients ($\alpha$, $\boldsymbol{\beta}$) are placeholder values
fit to synthetic data structured to match VCH's PARIS system; the framework is designed
for direct replacement with real logistic regression estimates once offer outcome data
are received.

For the BIP evaluation (Section~\ref{sec:bip_results}), a synthetic snapshot of
$P = 500$ clients and $N = 30$ facilities (total 3{,}344 beds) is drawn from the same
distributions.  The resulting $500 \times 30$ deferral probability matrix has mean
$\bar{p}_{pn} = 0.187$, consistent with VCH's baseline declination environment.

\subsubsection*{Reconciliation of Published Wait-Time Figures}

The OSA reports two wait-time statistics for VCH that appear contradictory: a mean wait
of 81 days across \emph{all new admissions} and 473 days for \emph{non-urgent community}
clients.  These refer to different populations and are fully consistent:

\begin{enumerate}[noitemsep]
  \item The 81-day figure covers all 147 VCH bed openings per month, including 129 direct
        hospital-to-LTC transfers that bypass the priority waitlist.
  \item The 473-day figure covers only the 18 waitlist-accessible vacancies per month
        that are within this study's scope.
\end{enumerate}

\noindent
Our model is calibrated to the 473-day target, which directly characterises the
priority waitlist.

\subsection{Results}
\label{sec:vchr}

We evaluate the allocation model at three levels.  The BIP gives the theoretical
minimum for a static snapshot.  The queue model evaluates dynamic online policies over
a 3-year simulation.  The DES adds escalation, batch arrivals, flu seasonality, and
geographic matching.

\subsubsection{BIP: Theoretical Benchmark}
\label{sec:bip_results}

The LP relaxation of the BIP (Section~\ref{sec:bip}) is solved on the 500-client
snapshot.  Because the assignment constraint matrix (C1 + C2) is the node-arc
incidence matrix of a bipartite graph, it is totally unimodular; the LP optimal
solution is therefore guaranteed integer without branch-and-bound.
Table~\ref{tab:bip_results} reports the four-policy comparison.

\begin{table}[H]
\centering
\caption{BIP policy comparison --- 500-client snapshot, 30 facilities, 3{,}344 beds.}
\label{tab:bip_results}
\begin{tabular}{lrrrr}
\toprule
\textbf{Metric} & \textbf{Historical} & \textbf{Cyclic FIFO} & \textbf{Random} & \textbf{LP Optimal} \\
\midrule
Clients assigned             & 467   & 500   & 500  & 500 \\
Total expected declinations  & 83.0  & 113.5 & 94.4 & 46.4 \\
Mean $\hat{p}_{pn}$          & 0.178 & 0.227 & 0.189 & 0.093 \\
Facilities used              & 30    & 5     & 30   & 5 \\
\bottomrule
\end{tabular}
\end{table}

\begin{table}[H]
\centering
\caption{BIP per-tier mean $\hat{p}_{pn}$ across policies.}
\label{tab:bip_tier}
\begin{tabular}{lrrrr}
\toprule
\textbf{Tier} & \textbf{Historical} & \textbf{Cyclic FIFO} & \textbf{Random} & \textbf{LP Optimal} \\
\midrule
Acute Care             & 0.154 & 0.189 & 0.147 & 0.083 \\
Community Emergency    & 0.176 & 0.211 & 0.187 & 0.096 \\
Community High         & 0.180 & 0.242 & 0.218 & 0.099 \\
Transfer/Site Specific & 0.250 & 0.342 & 0.251 & 0.100 \\
\bottomrule
\end{tabular}
\end{table}

\noindent
The LP optimal assignment reduces total expected declinations from 113.5 to 46.4
($-59.1\%$) relative to Cyclic FIFO, and from 94.4 to 46.4 ($-50.8\%$) relative to
random assignment.  The largest per-tier gain is for Transfer/Site Specific clients
($0.342 \to 0.100$, $-70.8\%$), who are hardest to place under Cyclic FIFO due to
site-specific facility constraints.  The LP achieves this not by distributing load
evenly (both policies concentrate clients in 5 facilities) but by identifying the
assignments with the lowest compatibility risk.

This 59\% figure is a theoretical upper bound: it assumes perfect information about all
$\hat{p}_{pn}$ values and a static waitlist.  The dynamic policies below show how much
of this gain is achievable in real-time operation.

\subsubsection{Queue Model: Online Simulation}
\label{sec:qm_results}

The queue model is run for a 3-year horizon (1{,}095 days) with a 6-month warm-up
(180 days), yielding 915 effective measurement days.  Thirty replications are used.

\begin{table}[H]
\centering
\caption{Queue model overall results (30 reps, 1{,}095-day horizon, 180-day warm-up).
         Values are mean $\pm$ 95\% CI accumulated over the effective period.}
\label{tab:vch_overall}
\begin{tabular}{lrrr}
\toprule
\textbf{Metric} & \textbf{Erlang-A} & \textbf{Cyclic FIFO} & \textbf{Optimised} \\
\midrule
Placements            & $542 \pm 7$    & $548 \pm 9$      & $531 \pm 7$ \\
Declinations          & $0 \pm 0$      & $118 \pm 5$      & $96 \pm 4$ \\
Abandonments          & $329 \pm 11$   & $309 \pm 11$     & $294 \pm 10$ \\
Bed-days wasted       & $0.0$          & $235.6 \pm 10.9$ & $191.7 \pm 7.3$ \\
Bed-days per placement & $0.000$       & $0.430 \pm 0.020$ & $0.360 \pm 0.010$ \\
Declination rate      & $0.000$        & $0.180 \pm 0.010$ & $0.150 \pm 0.000$ \\
Mean wait (days)      & $55.3 \pm 2.4$ & $68.8 \pm 4.2$   & $34.7 \pm 1.9$ \\
Mean queue length     & $44.8 \pm 1.9$ & $61.5 \pm 3.3$   & $89.6 \pm 4.0$ \\
\bottomrule
\end{tabular}
\end{table}

\begin{table}[H]
\centering
\caption{Queue model per-tier results (mean $\pm$ 95\% CI).
         EA = Erlang-A; CF = Cyclic FIFO; Opt = Optimised.}
\label{tab:vch_tier}
\begin{tabular}{lrrrrrr}
\toprule
& \multicolumn{3}{c}{\textbf{Declination rate}} & \multicolumn{3}{c}{\textbf{Mean wait (days)}} \\
\cmidrule(lr){2-4} \cmidrule(lr){5-7}
\textbf{Tier} & \textbf{EA} & \textbf{CF} & \textbf{Opt} & \textbf{EA} & \textbf{CF} & \textbf{Opt} \\
\midrule
Transfer/Site & 0.000 & $0.230 \pm 0.010$ & $0.200 \pm 0.020$ & $59.7 \pm 2.8$ & $36.4 \pm 3.7$ & $42.5 \pm 3.4$ \\
Comm.\ High   & 0.000 & $0.190 \pm 0.010$ & $0.170 \pm 0.010$ & $68.5 \pm 3.7$ & $146.9 \pm 11.8$ & $100.4 \pm 9.0$ \\
Comm.\ Emerg. & 0.000 & $0.160 \pm 0.010$ & $0.150 \pm 0.010$ & $49.6 \pm 2.4$ & $49.8 \pm 2.9$ & $21.5 \pm 1.5$ \\
Acute Care    & 0.000 & $0.140 \pm 0.010$ & $0.120 \pm 0.010$ & $40.6 \pm 2.5$ & $32.8 \pm 1.3$ & $9.5 \pm 0.5$ \\
\bottomrule
\end{tabular}
\end{table}

\noindent
The optimised policy reduces the declination rate from 0.180 to 0.150 ($-18.7\%$),
saves $235.6 - 191.7 \approx 44$ bed-days over the simulation period, and cuts Acute
Care mean wait from 32.8 to 9.5 days ($-71.0\%$).  Community High wait also improves
from 146.9 to 100.4 days ($-31.7\%$), showing that the fairness window prevents
starvation of non-urgent tiers while still directing vacancies to better-matched clients.
Transfer/Site wait increases modestly ($+16.7\%$), reflecting the reallocation of some
vacancies to tiers with lower deferral risk.

The Erlang-A policy provides a useful reference: it eliminates all declinations but uses
pure FCFS across tiers, yielding a \emph{higher} Acute Care wait (40.6 days) than Cyclic
FIFO (32.8 days).  Cyclic FIFO's equal-rotation structure guarantees Acute Care clients
regular service slots regardless of queue depth, whereas FCFS does not.  This confirms
that zero declinations and short urgent-tier waits are distinct objectives.

\subsubsection{DES: Extended Simulation}
\label{sec:des_results}

The DES extends the queue model with four additions: (i)~priority escalation
(CommEm $\to$ AC after 45 days; Transfer $\to$ CommHigh after 90 days);
(ii)~batch arrivals (10\% of arrival events bring 2--4 clients, modelling
hospital discharge planning days); (iii)~flu-season deferral adjustment
(November--January); and (iv)~a geographic soft-penalty $\gamma = 0.25$ added
to the scoring function, reflecting client and family preference for in-region placement.

\begin{table}[H]
\centering
\caption{DES overall results (30 reps, 1{,}095-day horizon, 180-day warm-up).}
\label{tab:des_overall}
\begin{tabular}{lrrr}
\toprule
\textbf{Metric} & \textbf{Erlang-A} & \textbf{Cyclic FIFO} & \textbf{Optimised} \\
\midrule
Placements           & $555 \pm 9$    & $550 \pm 10$    & $558 \pm 7$ \\
Declinations         & $0 \pm 0$      & $184 \pm 7$     & $138 \pm 5$ \\
Abandonments         & $470 \pm 15$   & $441 \pm 16$    & $394 \pm 14$ \\
Priority escalations & $210 \pm 9$    & $63 \pm 4$      & $149 \pm 6$ \\
Bed-days wasted      & $0.0$          & $368.0 \pm 13.7$ & $276.1 \pm 10.7$ \\
Bed-days/placement   & $0.000$        & $0.670 \pm 0.020$ & $0.490 \pm 0.020$ \\
Declination rate     & $0.000$        & $0.250 \pm 0.010$ & $0.200 \pm 0.010$ \\
Mean wait (days)     & $75.5 \pm 3.0$ & $106.1 \pm 5.0$  & $34.7 \pm 1.3$ \\
Flu-season declin.\ rate & $0.000$    & $0.280 \pm 0.020$ & $0.210 \pm 0.010$ \\
Out-of-region fraction & $0.750 \pm 0.010$ & $0.730 \pm 0.010$ & $0.290 \pm 0.010$ \\
\bottomrule
\end{tabular}
\end{table}

\begin{table}[H]
\centering
\caption{DES per-tier results (mean $\pm$ 95\% CI).
         EA = Erlang-A; CF = Cyclic FIFO; Opt = Optimised.}
\label{tab:des_tier}
\begin{tabular}{lrrrrrr}
\toprule
& \multicolumn{3}{c}{\textbf{Declination rate}} & \multicolumn{3}{c}{\textbf{Mean wait (days)}} \\
\cmidrule(lr){2-4} \cmidrule(lr){5-7}
\textbf{Tier} & \textbf{EA} & \textbf{CF} & \textbf{Opt} & \textbf{EA} & \textbf{CF} & \textbf{Opt} \\
\midrule
Transfer/Site & 0.000 & $0.300 \pm 0.010$ & $0.260 \pm 0.010$ & $70.4 \pm 1.9$ & $26.8 \pm 3.4$ & $22.9 \pm 0.9$ \\
Comm.\ High   & 0.000 & $0.260 \pm 0.010$ & $0.230 \pm 0.010$ & $78.2 \pm 3.3$ & $272.7 \pm 16.6$ & $106.1 \pm 8.4$ \\
Comm.\ Emerg. & 0.000 & $0.240 \pm 0.010$ & $0.190 \pm 0.010$ & $42.0 \pm 0.9$ & $36.6 \pm 0.5$ & $11.5 \pm 0.4$ \\
Acute Care    & 0.000 & $0.200 \pm 0.010$ & $0.170 \pm 0.010$ & $74.8 \pm 2.9$ & $78.9 \pm 3.4$ & $25.5 \pm 0.9$ \\
\bottomrule
\end{tabular}
\end{table}

Key DES findings:
\begin{itemize}[noitemsep]
  \item \textbf{Declinations:} reduced from 0.250 to 0.200 ($-20.0\%$); the flu-season
        rate drops from 0.280 to 0.210 ($-25.0\%$), showing that the optimised policy
        adapts to seasonal fluctuations without explicit seasonal targeting.
  \item \textbf{Waiting times:} overall mean falls from 106 to 35 days ($-67.3\%$),
        with the largest reductions for Community Emergency ($-68.6\%$) and Acute Care
        ($-67.7\%$).
  \item \textbf{Geographic matching:} out-of-region fraction falls from 73\% to 29\%
        (in-region placements: 27\% $\to$ 71\%).  The geographic soft-penalty in the
        optimised scoring function naturally favours in-region facilities when they
        offer comparable deferral risk, without any explicit geographic quota.
  \item \textbf{Priority escalations:} increase from 63 to 149 under the optimised policy.
        Under Cyclic FIFO, the steady 25\% rotation ensures most CommEm clients are placed
        before the 45-day escalation threshold.  Under the optimised policy, clients with
        higher $\hat{p}_{pn}$ are selected less frequently in early service windows; some
        wait past the escalation threshold before being matched.  This is a fairness
        consideration to monitor in practice --- escalating clients incur the cost of a
        priority change, which may affect care planning.
\end{itemize}

\subsubsection{Summary Across All Three Models}
\label{sec:results_summary}

Table~\ref{tab:summary_all} consolidates the key improvements across BIP, queue model,
and DES.

\begin{table}[H]
\centering
\caption{Summary of improvement (Cyclic FIFO $\to$ LP Optimal / Optimised) across models.}
\label{tab:summary_all}
\begin{tabular}{llrrc}
\toprule
\textbf{Model} & \textbf{Metric} & \textbf{Cyclic FIFO} & \textbf{Optimised / LP} & \textbf{Change} \\
\midrule
BIP (static)   & Expected declinations        & 113.5 & 46.4  & $-59.1\%$ \\
               & Mean $\hat{p}_{pn}$          & 0.227 & 0.093 & $-59.0\%$ \\
\midrule
Queue Model    & Declination rate             & 0.180 & 0.150 & $-18.7\%$ \\
               & Acute Care wait (days)       & 32.8  & 9.5   & $-71.0\%$ \\
               & Community High wait (days)   & 146.9 & 100.4 & $-31.7\%$ \\
               & Bed-days wasted              & 235.6 & 191.7 & $-18.6\%$ \\
\midrule
DES            & Declination rate             & 0.250 & 0.200 & $-20.0\%$ \\
               & Overall mean wait (days)     & 106.1 & 34.7  & $-67.3\%$ \\
               & Bed-days wasted              & 368.0 & 276.1 & $-25.0\%$ \\
               & In-region placement fraction & 27\%  & 71\%  & $+163\%$ \\
\bottomrule
\end{tabular}
\end{table}

\noindent
The gap between the BIP bound ($-59.1\%$) and the online policies ($-18.7\%$ to
$-20.0\%$) reflects the cost of operating with incomplete future information in a
dynamic queue.  The BIP solves globally over a known, static waitlist; the online
policies observe only the current queue state and must balance efficiency across tiers.
Nonetheless, the fairness-window policy captures a meaningful fraction of the
theoretical gain and delivers consistent improvements on all metrics --- declinations,
wait times, bed utilisation, and geographic matching --- without any additional
system resources.
